% Chapter 2 -- Rubric: PO1 (1.1), PO2 (research literature), feeds PO3 and PO11
\chapter{Preliminary Research}
\label{ch:research}

\section{Literature Review}
\label{sec:literature}

Eleven primary works were read in full and form the evidentiary basis of this
review. Seven are surveyed comparatively in \Cref{tab:llm-systems}: four GUI
\emph{testing} systems~\cite{gptdroid,droidagent,auitestagent,xuatcopilot} and
three GUI \emph{task-automation} agents~\cite{autodroid,appagent,appagentx}
whose memory and grounding mechanisms transfer directly to the testing problem.
Four further works are treated separately in \Cref{subsec:formative} because
they did not merely inform this review but shaped the architecture that
\Cref{ch:design} presents. Every figure quoted below was taken from those
papers. The review is organised around five questions: how the field arrived at
language-model agents, which works directly formed this design, what those
agents demonstrably add, how the resulting claims are measured, and what the
literature leaves open.

\subsection{From Random Input to Learned Exploration}
\label{subsec:lineage}

Automated Android GUI testing has passed through three generations. The first
drives the application with pseudo-random events; Android's own Monkey is the
canonical example and remains the universal baseline~\cite{monkey}. The second
is systematic or model-based: DroidBot builds a UI transition graph and
greedily prioritises unexplored widgets~\cite{droidbot}. The third is
learning-based: Humanoid trains a deep network on real human interaction traces
to bias action selection toward human-like paths~\cite{humanoid}, and a
parallel line applies reinforcement learning to reach novel states. Q-testing
makes the reward explicitly curiosity-driven, keeping a memory of visited
states and training a neural state-comparison module to judge state novelty at
the granularity of functional scenarios rather than of widgets; on 50
applications it exceeded the coverage and fault detection of the tools of its
day, and 22 of its reported faults were confirmed~\cite{qtesting}. ARES
replaces tabular reinforcement learning with deep networks for black-box
Android testing, on the argument that the state space of a real application is
far too large for a fixed table~\cite{ares}.

Liu et al.\ summarise the limitations of the learning-based generation
precisely~\cite{gptdroid}: these methods require large interaction datasets
that are difficult to collect; they generalise poorly to applications they were
not trained on, while applications evolve continuously; and the environment is
non-deterministic. Their own example is that deleting the last item from a list
produces an empty list on which the delete control no longer works, which
defeats a learned policy.

A second, more fundamental criticism runs through the whole lineage and is the
one this project takes as its starting point. All three generations optimise
\emph{structural} objectives: activity coverage, widget coverage, code coverage.
Yoon et al.\ point out that this is not what practitioners
want~\cite{droidagent}: an empirical study of Android developers found that the
dominant test-design strategy is to follow the usage model of the application,
that developers prefer test cases targeting individual features and use cases,
and that, asked what form they would like generated tests to take, they chose
natural language over test-API scripts, with expected outputs and use-case
steps included. Coverage is a proxy that the people who consume the results did
not ask for~\cite{coppola}.

\subsection{The Language-Model Agent Wave}
\label{subsec:llmwave}

Between 2023 and 2025 the field reorganised around large-language-model agents.
\Cref{tab:llm-systems} summarises the seven systems reviewed.

\begingroup
\footnotesize
\setlength{\tabcolsep}{4pt}
\begin{xltabular}{\textwidth}{@{}>{\hsize=0.52\hsize}Y>{\hsize=1.02\hsize}Y>{\hsize=1.46\hsize}Y@{}}
  \caption{The seven primary systems reviewed: goal, core mechanism, and the basis on which each was evaluated}
  \label{tab:llm-systems}\\
  \toprule
  System & Goal and mechanism & Evaluation basis and headline result \\
  \midrule
  \endfirsthead
  \caption[]{The seven primary systems reviewed \emph{(continued)}}\\
  \toprule
  System & Goal and mechanism & Evaluation basis and headline result \\
  \midrule
  \endhead
  \midrule
  \multicolumn{3}{r@{}}{\footnotesize\itshape continued on the next page}\\
  \endfoot
  \bottomrule
  \endlastfoot

  GPTDroid \cite{gptdroid} & GUI \emph{testing} framed as a question-and-answer task; functionality-aware memory prompting carries testing knowledge across the whole session. &
    93 Android applications carrying 143 previously confirmed crash bugs. 75\% activity and 66\% code coverage against 57\% for the best baseline; 95 of the known bugs detected, 31\% more than the best baseline, and 35\% faster. Separately, 53 previously unknown crashes on Google Play, of which 35 were confirmed or fixed. \\
  \addlinespace
  DroidAgent \cite{droidagent} & GUI \emph{testing} with four cooperating agents (Planner, Actor, Observer, Reflector) over short-term, long-term and spatial memory; sets its own natural-language task goals. &
    15 applications from the Themis benchmark \cite{themis}, two-hour budget each. 61\% activity coverage against 51\% for Humanoid; 13.9 features covered per application against 7.3. Of 374 self-generated tasks, 85\% judged viable and 59\% completed by manual labelling. Five crashes found. Cost \$13--\$22 per application, mean \$18.10. \\
  \addlinespace
  AUITestAgent \cite{auitestagent} & GUI \emph{function testing} driven by a natural-language test requirement; deliberately decouples interaction from verification and extracts verification evidence along three dimensions of the interaction log. &
    30 authored interaction tasks over 8 popular applications; 77\% completion. Verification measured on 40 tasks built from 20 \emph{manually injected} anomalies: 95\% oracle accuracy on correct behaviour, 85\% on anomaly detection, 90\% recall at a 4.5\% false-positive rate. Field-deployed at Meituan: 4 new functional bugs across 10 regression rounds in two months, saving 1.5 person-days per round. \\
  \addlinespace
  XUAT-Copilot \cite{xuatcopilot} & Automating the \emph{test-script generation} stage of an existing industrial user-acceptance-testing pipeline; three agents for action planning, state checking and parameter selection, plus reflection. &
    450 WeChat Pay test cases, all of which had already passed with human-written scripts. $Pass@1$ of 88.55\% and $Complete@1$ of 93.03\%, against 22.65\% and 25.25\% for a single-agent variant carrying the same combined prompt. \\
  \addlinespace
  AutoDroid \cite{autodroid} & \emph{Task automation}: combines model commonsense with app-specific knowledge mined by automated dynamic analysis; UI transition memory, memory injection, and query optimisation. &
    DroidTask benchmark: 158 tasks over 13 applications. With GPT-4, 90.9\% action accuracy and 71.3\% task completion, 36--40 percentage points above the language-model baselines. Query cost reduced from \$18.76 to \$10.17 per thousand GPT-4 queries by halving prompt length. \\
  \addlinespace
  AppAgent \cite{appagent} & \emph{Task automation} from screenshots alone; an exploration phase writes a knowledge-base document per application, either autonomously or by watching a human demonstration. &
    45 tasks over 9 applications. Success rate 2.2\% with a raw action space, 48.9\% with a designed action space and no document, 73.3\% with an autonomously explored document, 84.4\% from watching demonstrations, and 95.6\% from a hand-written document. \\
  \addlinespace
  AppAgentX \cite{appagentx} & \emph{Task automation} with an evolutionary memory: repeated low-level action sequences are abstracted into high-level shortcuts, implemented over Neo4j for graph storage and a vector index for retrieval. &
    AppAgent benchmark (50 tasks, 9 applications), each repeated five times. Success rate 16.9\% with no memory, 70.8\% with chain-structured memory, 71.4\% with memory plus evolution, while average steps fell from 9.1 to 5.7 and tokens per task from 9.26k to 4.94k. \\
\end{xltabular}
\endgroup

\subsection{Directly Formative Work}
\label{subsec:formative}

Four works are separated from \Cref{tab:llm-systems} because they did not
merely inform this review: they supplied the architecture, the knowledge
representation, the feedback loop and the execution baseline that this project
adopted. They are reported here with the evidence that justified adopting them.

\paragraph{The multi-agent decomposition (the architectural source).} Su et
al.'s work on automated soap opera testing is the direct motivation for this
project's multi-agent structure~\cite{soapopera}. Soap opera testing is
scenario-based exploratory testing: complex, narrative test scenarios designed
to provoke failures, which a tester executes while remaining alert for anything
unexpected along the way. Their system decomposes the problem into three
multi-turn-dialogue agents over a Scenario Knowledge Graph (SKG). The
\emph{Planner} determines the next step and expands it into executable
sub-steps, retrieving step knowledge from the SKG; the \emph{Player} translates
one sub-step into a structured UI instruction and executes it; and the
\emph{Detector} runs \emph{after every Player execution}, retrieving the
oracles associated with that step and judging whether the resulting interface
state violates any of them.

That division maps almost directly onto the planner, executor and investigator
roles of \Cref{ch:design}, and the mapping was deliberate. Two details of their
design were adopted and one was deliberately not. Adopted: the separation of
the agent that acts from the agent that judges, and the retrieval of
\emph{oracles} rather than merely steps, so that judgement is grounded in
recorded expectations rather than in the model's opinion. Not adopted: their
Detector runs after every single step, which is the stronger design and is
recorded in \Cref{ch:conclusion} as outstanding work; this project's
investigator is invoked at verdict time.

\paragraph{Why that decomposition is not sufficient on its own.} The same paper
supplies the most direct quantitative warning available to this project, and it
is the reason the investigator role was strengthened rather than merely
included. Over 40 soap opera tests on four applications, their automated system
reported \emph{more} bugs than the human participants (158 against 80) while
being right far less often: true-positive accuracy of 0.516, 0.656, 0.577 and
0.605 across the four applications, against 1.000 for the manual
baseline~\cite{soapopera}. Roughly two in five reported defects were false. Of
80 bugs found manually and 158 found automatically, only six were the same
bug.

Their analysis of the 67 incorrect reports is directly actionable and informed
this project's reporting rules: 41.8\% were unreasonable or unnecessary
\emph{improvement suggestions} rather than defects, 29.9\% were model
hallucinations (including bugs asserted about execution paths that were never
taken), 17.9\% came from screenshots captured before rendering completed, and
10.4\% from mislocating the element that was acted on. Their ablation is
equally pointed: removing oracle knowledge from the Detector roughly halves
detection accuracy, from 0.516 to 0.216 on Firefox and from 0.656 to 0.421 on
WordPress. An agent that judges without recorded expectations does not fail
gracefully; it becomes confidently wrong. Their reported cost, approximately
US\$1.76 per scenario and US\$0.77 per true-positive defect, is also one of the
few per-defect cost figures in the literature.

\paragraph{Knowledge graphs for exploratory testing (the knowledge
representation).} The predecessor from the same group establishes that a
knowledge graph built by language models is a workable substrate for
exploratory testing~\cite{etkg}. They construct a System Knowledge Graph of
User Tasks and Failures from 32,772 Mozilla Firefox bug reports by extracting
preconditions, steps to reproduce, expected behaviours and observed behaviours,
splitting compound steps into atomic single-operation steps, and clustering
equivalent steps so that scenarios connect through shared nodes rather than
standing alone. Chain-of-thought reasoning then filters generated scenarios for
feasibility, which rule-based predecessors could not do.

Two findings carried into this project's design. First, on data quality: only
13,943 of the 32,772 bug reports contained clearly delineated sections, so more
than half the corpus was unusable by the rule-based extraction that preceded
them; the tolerance of language-model extraction to badly structured input is
what makes graph construction from real project data feasible at all. Second,
on usefulness: in a study of 16 participants, the eight using their generated
scenarios found 36 bugs in two hours against 21 for the state-of-the-art
baseline. Note the shape of that evaluation, however, because it marks the
boundary of the comparison: their system generates scenarios \emph{for human
testers to execute}, whereas this project executes autonomously and must
therefore supply its own oracle.

\paragraph{The feedback loop has a pre-LLM precedent.} It is worth recording
that the central idea of this project, a specification-derived model corrected
by what execution actually observes, is not new with language models. Gebizli
and S\"ozer's ARME refines model-based-testing models using exploratory testing
traces: because such models are built by hand from functional requirements they
omit real usage behaviour, so recorded exploratory traces are compared against
the model's paths, omitted behaviour is added, and Markov-chain transition
probabilities are updated by observed visit frequency, iteratively~\cite{arme}.
Across three industrial case studies on a digital television system the refined
models yielded 5, 7 and 6 new faults, including one that reset the device and
one that crashed it; critically, those faults were found by neither the
original test suite nor the exploratory sessions that produced the traces.

This project's contribution is therefore not the idea of closing the loop
between an a-priori specification and a-posteriori observation, which ARME
demonstrated in 2017, but its instantiation: a graph rather than a Markov
chain, requirements ingested automatically rather than modelled by hand,
exploration performed by an agent rather than by the engineers whose traces
ARME consumed, and defect attribution as a first-class outcome.

\paragraph{Natural-language execution, and the closest institutional
comparator.} LELANTE executes a natural-language test case description directly
on Android without a pre-written script, through screen refinement, a
structured prompt and chain-of-thought action generation, with model
distillation to control cost~\cite{lelante}. Across 390 test cases on ten
applications it reports a 73\% test-execution success rate.

It is the closest comparator available for this project's executor considered
alone, and the comparison is bounded in a way worth stating: LELANTE executes a
test case it is given, and neither decides what to test nor judges whether the
application behaved correctly. Its 73\% is therefore a figure for the
actuation problem only. It is also the work closest to this project
institutionally, sharing both the university and the industry partner, which
makes its reported limitations a reasonable prior for what this project should
expect on the same class of application.

\subsection{What the Evidence Agrees On}
\label{subsec:agreement}

Three findings recur across systems that were built independently, by different
groups, on different benchmarks. Because they agree, they can be treated as
established rather than as any one paper's claim.

\paragraph{Persistent memory is the single largest lever, and it is worth more
than the model.} The ablations are unusually consistent. AppAgentX reports
16.9\% success with no memory and 70.8\% with chain-structured memory, a
four-fold improvement from memory alone, before any other
enhancement~\cite{appagentx}. AppAgent reports 2.2\% with a raw action space
and no document, rising to 73.3\% once an exploration phase has written one, on
the same underlying model~\cite{appagent}. Removing the explored-functionality
memory from GPTDroid costs it 51\% of its activity coverage, from 0.75 to 0.37,
the largest single contribution of any of its modules~\cite{gptdroid}.
AutoDroid's entire thesis is that app-specific knowledge mined by dynamic
analysis, injected into the prompt, is what separates 71.3\% task completion
from the 31--35\% of the unaugmented baselines~\cite{autodroid}. This is the
strongest single justification available for the knowledge-graph architecture
selected in \Cref{sec:alternatives}.

\paragraph{Decomposing the agent into roles beats one large prompt.} XUAT-Copilot
gives the cleanest measurement, because its single-agent variant is constructed
by concatenating all three agents' prompts into one: $Pass@1$ falls from
88.55\% to 22.65\%~\cite{xuatcopilot}. The authors attribute the collapse to
long prompts inducing forgetting, hallucination and inefficient planning.
DroidAgent reaches the same architecture independently with four
roles~\cite{droidagent}, AUITestAgent makes the separation of interaction from
verification its central design claim~\cite{auitestagent}, and the soap opera
system of \Cref{subsec:formative} arrives at a three-role Planner, Player and
Detector split for the specifically exploratory case~\cite{soapopera}. Four
groups converging on role decomposition from four different starting points is
the strongest form this review's evidence takes.

\paragraph{The bottleneck is the oracle, not the actuator.} Across the seven
systems, \emph{interacting} with an application is largely solved: completion
rates cluster between 71\% and 95\%. \emph{Judging} the result is not. The
sharpest evidence is AUITestAgent's baseline, a competent general-purpose model
prompted directly: it reaches 0.90 oracle accuracy when the application is
behaving correctly and 0.30 when an anomaly is present, and the authors name
the cause: ``its tendency to output a `no anomaly'
judgment''~\cite{auitestagent}. A verifier that is accurate only when there is
nothing to find is not a verifier; it is a silent pass. This single observation
motivated two design decisions in this project: the separate investigator role
(DR-06) and the prohibition on silent fallbacks (NFR-D2).

The soap opera results of \Cref{subsec:formative} show the same bottleneck from
the opposite direction and are worth setting beside it, because the two failure
modes are mirror images. AUITestAgent's baseline errs toward silence, declaring
nothing wrong when something is; the soap opera Detector errs toward noise,
reporting 158 defects at 0.516--0.656 accuracy where humans reported 80 at
1.000, with 41.8\% of its errors being improvement suggestions rather than
defects at all~\cite{soapopera}. Both are the same underlying failure, an
oracle unanchored to a recorded expectation, and together they bound the design
problem: a useful verifier must be able to say \emph{nothing is wrong here}
without that becoming its default, and must distinguish a defect from a
preference. This is why this project's investigator is required to cite the
observation supporting its verdict, and why a verdict unsupported by evidence
is recorded as unverified rather than resolved in either direction.

A fourth observation is weaker but practically important: \textbf{cost is
rarely reported, and where it is reported it is substantial.} DroidAgent spends
a mean of \$18.10 per application for a two-hour budget and states in its
threats to validity that ``given the monetary constraints linked to API
requests, we could not conduct multiple runs''~\cite{droidagent}. That is, the
cost of the method prevented the authors from establishing the reproducibility
of their own result. AutoDroid treats query cost as a first-class engineering
target and halves it by compressing the UI representation~\cite{autodroid}. Of
the remaining five systems, only AUITestAgent and AppAgentX report token
consumption at all.

\subsection{How This Literature Is Measured, and Why It Reports Almost No Precision or Recall}
\label{subsec:measurement}

The most consequential finding of this review is methodological. Across all
seven systems, defect-finding ability is established in one of three ways, none
of which supports a precision, recall or $F_1$ figure.
\Cref{tab:measurement} sets out what each regime can and cannot establish.

\begingroup
\small
\setlength{\tabcolsep}{5pt}
\begin{xltabular}{\textwidth}{@{}>{\hsize=0.62\hsize}Y>{\hsize=1.08\hsize}Y>{\hsize=1.30\hsize}Y@{}}
  \caption{The three evaluation regimes in the reviewed literature, and the limit of each}
  \label{tab:measurement}\\
  \toprule
  Regime & Used by & What it can and cannot establish \\
  \midrule
  \endfirsthead
  \caption[]{The three evaluation regimes \emph{(continued)}}\\
  \toprule
  Regime & Used by & What it can and cannot establish \\
  \midrule
  \endhead
  \midrule
  \multicolumn{3}{r@{}}{\footnotesize\itshape continued on the next page}\\
  \endfoot
  \bottomrule
  \endlastfoot

  Structural coverage proxies (activity, widget, code, feature) &
    GPTDroid, DroidAgent, and every pre-LLM baseline &
    Establishes that more of the application was reached. Cannot establish that anything was \emph{checked}: an activity can be covered by a control that opens it and a back press, which DroidAgent documents Humanoid doing. Feature coverage is stronger but, absent a specification, must be defined by hand: DroidAgent's feature list was fixed ``until the consensus of three authors''. \\
  \addlinespace
  Known-crash benchmarks &
    GPTDroid on Themis \cite{themis} plus 73 further applications: 93 applications, 143 bugs &
    Gives a genuine denominator and therefore a genuine recall, 95 of 143. Its limit is stated by the authors themselves: ``same as the benchmark, all bugs are crash bugs.'' A crash is the only defect a tool can recognise without knowing what the application was supposed to do, so the entire quantitative bug-finding literature measures crash detection. Functional defects (a wrong total, an order routed to the wrong seller) are outside the benchmark by construction. \\
  \addlinespace
  Wild-bug counts &
    GPTDroid (53 new, 35 confirmed or fixed); AUITestAgent (4 in two months at Meituan) &
    Strong evidence of practical usefulness, and a precision of sorts can be computed from the confirmation rate. Recall is unobtainable in principle: nobody knows how many defects a production application contains, so there is no denominator. \\
\end{xltabular}
\endgroup

Only two of the seven papers report precision and recall at all, and it is
instructive that both had to manufacture their ground truth to do so.

AUITestAgent needed to measure verification ability on commercial applications,
and states the obstacle plainly: ``as the apps we selected are not open-source
and bugs are not common in the online versions of these mature commercial apps,
we needed a robust method to effectively measure AUITestAgent's performance in
executing function verification''~\cite{auitestagent}. Their solution was to
\emph{inject} defects: two authors analysed 50 recent historical anomalies at
Meituan, distilled them into three categories, and manually created 20
corresponding anomaly manifestations, each paired with one correct and one
anomalous interaction trace. Only against that artificial ground truth could
they report 90\% recall at a 4.5\% false-positive rate.

DroidAgent reports precision 0.72, recall 0.77 and $F_1$ 0.74, but for its
\emph{Reflector's self-assessment}, not for defect detection~\cite{droidagent}.
The measurement is of whether the agent correctly judged that it had completed
its own task, scored against three authors' manual labels. It is a measure of
self-knowledge, which is valuable and which this project also needs, but it
says nothing about how many defects were missed.

A fourth regime appears in the formative work and is the most honest of the
four, though it scales poorly. The soap opera study establishes validity by
having two authors independently reproduce every reported defect on a real
device, resolve disagreements by discussion, and then submit the surviving
reports to the projects themselves for developer
adjudication~\cite{soapopera}. This yields a genuine \emph{precision} (their
0.516--0.656 accuracy figures) because each report is checked individually, and
it is why that study can state how wrong its agent is where the seven cannot.
It still yields no recall, for the same reason wild-bug counts cannot: there is
no denominator. The cost is human effort proportional to the number of reports,
which is precisely what automation was meant to avoid, and it is the regime
this project's own campaign reports approximate when a human confirms a finding.

The pattern is therefore a chain of dependencies, and it is worth stating
explicitly because it determines what this project can and cannot claim:
\emph{without a specification there is no oracle beyond a crash; without an
oracle there is no labelled ground truth; without ground truth there is no
precision, recall or $F_1$, only coverage proxies, crash counts, and human
consensus.} Every subjective judgement in this literature (DroidAgent's three
authors agreeing on a feature list, its manual labelling of 374 tasks, the
confirmation of 35 of 53 reported Google Play bugs by their developers) exists
to substitute for a ground truth that the absence of a specification made
unavailable.

This diagnosis has two direct consequences for the present project. First, it
is the strongest argument for the specification-grounded design: ingesting a
requirements document supplies exactly the oracle whose absence forces everyone
else onto coverage proxies, and makes \emph{requirement coverage} reportable as
a first-class result rather than activity coverage as a proxy. Second, it means
that this project's own precision and recall cannot be obtained honestly
without seeding known defects, in the manner of AUITestAgent. That experiment
is specified but was not completed within this session, because the industry
partner supplied neither a defect export nor a seeded build (constraint C3),
and it is carried into \Cref{ch:conclusion} as the single most valuable piece
of outstanding work.

\subsection{What the Literature Leaves Open}
\label{subsec:gaps}

\begin{enumerate}[label=\textbf{G\arabic*.},leftmargin=3.0em]
  \item \textbf{No specification is ever ingested.} Not one of the eleven works
    reads a requirements document, builds a requirement graph, or reports
    requirement coverage. Two come close from different directions and neither
    closes it. AUITestAgent accepts one natural-language test requirement per
    test, but holds no persistent model of the requirement set and cannot say
    which requirements remain untested~\cite{auitestagent}. The soap opera
    system holds exactly the persistent oracle model that AUITestAgent lacks,
    but derives it from bug reports and user manuals~\cite{soapopera,etkg}:
    knowledge of how the software \emph{has failed} and how people describe
    using it, not of what it was \emph{specified} to do. The distinction
    matters for coverage, because a bug-derived graph has no denominator. It
    can say a scenario was exercised; it cannot say which requirements remain
    untested, because it never held the requirement set.
  \item \textbf{Memory does not leave the application or the session.}
    AppAgentX's shortcut chains, AutoDroid's transition memory and GPTDroid's
    functionality memory are all per-application. No reviewed system transfers
    what it learned about one application, device profile or platform to
    another.
  \item \textbf{Defect history is rarely a knowledge source, and where it is,
    it is used for oracles rather than for prioritisation.} None of the seven
    systems in \Cref{tab:llm-systems} reads a bug tracker. The formative works
    of \Cref{subsec:formative} are the exception and must be stated as such:
    both build their knowledge graph \emph{entirely} from bug
    reports~\cite{etkg,soapopera}, which is the strongest existing evidence
    that defect history is a usable knowledge source at all. What remains open
    is the second half of the claim. Their bug-derived knowledge supplies
    scenarios and oracles, that is, \emph{what to check}; neither uses defect
    density to decide \emph{where to look first}, despite that being among the
    oldest and best-established predictors in software engineering. This
    project's defect source (ETA-REQ-301) is therefore a narrower contribution
    than a first reading of the seven systems would suggest, and it is
    positioned accordingly in \Cref{ch:design}.
  \item \textbf{Failures are not attributed.} Outcomes are binary. Nothing in
    these systems distinguishes ``the application is wrong'' from ``the agent
    could not drive the application'' from ``the environment blocked the test'',
    even though the three demand completely different responses and only the
    first is a defect report.
  \item \textbf{Web and mobile are studied in isolation.} The mobile systems
    reviewed here share no substrate with their web equivalents, so nothing
    learned on one platform informs the other.
  \item \textbf{Safety is largely unaddressed, and its cost is unmeasured.}
    Only AutoDroid treats destructive actions as a design problem, introducing
    privacy replacement and security confirmation and reporting what they cost:
    action accuracy falls from 92.9\% to 89.9\% and completion from 75.4\% to
    69.9\%~\cite{autodroid}. That roughly 3--6 percentage-point safety tax is the
    only quantification of this trade-off found in the review, and it is directly
    relevant to a project constrained to avoid irreversible actions on a
    production system.
\end{enumerate}

Gaps G1--G4 correspond, respectively, to this project's specification-grounded
planner, its multi-dimensional graph partitioning (ETA-REQ-304), its
defect-history source (ETA-REQ-301) and its three-way attribution taxonomy
(NFR-D3); G5 corresponds to the shared planner and graph behind the Android and
web executors (DR-04); G6 to the enforced safety boundary (DR-07). The
combination (a specification graph that says what \emph{should} happen, a
runtime graph that records what \emph{does}, and exploration steered by the gap
between them) is the position this project occupies in the literature.

That position should be claimed precisely, because one half of it has clear
prior art. Correcting a specification-derived model against what execution
observes is ARME's contribution from 2017, and it was demonstrated to find
faults that neither the original suite nor the exploratory sessions
found~\cite{arme}; the idea is not new here. What is new is the combination of
that loop with the three properties no reviewed system holds together: the
model is a requirement graph ingested automatically rather than a hand-built
Markov chain, the exploration that corrects it is performed by an autonomous
agent rather than by engineers whose traces are harvested afterwards, and every
outcome is attributed to the application, the agent or the environment before
it is allowed to change either graph. Stated at the level of the loop alone,
this project follows ARME; stated at the level of what fills the loop, it does
not.

\subsection{Additional Works Surveyed}
\label{subsec:additional}

Beyond the eleven works read in full, the team surveyed a wider set of recent
work on generalist GUI agents, visual grounding models and agent benchmarks.
Two strands bear directly on the argument of this chapter and are worth
reporting; the remainder are noted only briefly.

\paragraph{Benchmark environments with programmatic oracles.} The measurement
problem of \Cref{subsec:measurement} has a partial answer that none of the
seven testing systems uses: define success by inspecting durable application or
system state after the fact, rather than by asking the agent, or a language
model reading the agent's transcript, whether it succeeded. AndroidWorld builds
116 programmatically checkable tasks across 20 real Android applications, each
with an initialisation, a state-based success check and a teardown, and
parameterises them dynamically so that a memorised phrasing does not
suffice~\cite{androidworld}. AppWorld goes further for the class of behaviours
that touch stored data: its tasks are graded by database-state unit tests,
which accept any valid solution path while also detecting \emph{unintended
collateral changes}, something a transcript-reading judge cannot see at
all~\cite{appworld}.

Two consequences follow. First, this is a strictly stronger oracle than
self-report for the subset of behaviours whose effects are externally
observable, and adopting it where possible is recorded as future work in
\Cref{ch:conclusion}. Second, and more sobering, the scores these benchmarks
report are low: the best baseline in AndroidWorld completes 30.6\% of
tasks~\cite{androidworld}, and the strongest model in AppWorld solves about
49\% of normal tasks and 30\% of challenge tasks~\cite{appworld}. Set beside
the 71\% to 95\% completion rates that self-graded systems report
(\Cref{tab:llm-systems}), the gap is itself evidence for this chapter's central
claim: much of the apparent competence of GUI agents is an artefact of letting
them mark their own work. A project whose output is a correctness judgement
cannot afford to ignore that.

\paragraph{Learning from interactive outcomes.} DigiRL argues that static
demonstrations do not capture the stochastic behaviour encountered while
controlling a real device, and trains a device-control policy by offline
followed by offline-to-online reinforcement learning against a
vision-language-model evaluator, raising a 1.5B-parameter model from 17.7\% to
67.2\% success without additional human demonstrations~\cite{digirl}. It is the
clearest external support for this project's premise that execution outcomes
must change future behaviour, and it marks the boundary of the approach taken
here: this system adapts through graph memory, strategy scores and prompt
construction rather than by updating model weights, which keeps it applicable
to hosted models it does not control but forgoes the gains a learned policy can
reach.

\paragraph{The remainder.} The survey also covered visual grounding and screenshot-only action models. These are nearer to this project than to most of the systems in \Cref{tab:llm-systems}, because this system already combines both signals: the executor and the planner each receive a screenshot alongside the structural description, and a perceptual hash resolves screen identity where the accessibility tree is uninformative, which is the situation constraint C4 describes. The open question they raise is not whether to use pixels but when, since an image carries a cost at every step it is attached to. The survey also covered one recent review of operating-system agents, useful for terminology and for positioning this work as a mobile agent with explicit \emph{testing} responsibilities rather than a general phone assistant.
None of these changes the gap analysis of \Cref{subsec:gaps}.

\paragraph{How these three were read.} The benchmark works cited in this
subsection were surveyed from their published abstracts, result tables and
metadata rather than read in full, and the reported figures are quoted from
those sources. They are used here only for the orders of magnitude they supply
and for the existence of programmatic oracles, neither of which turns on detail
beyond the abstract. The seven testing systems of \Cref{tab:llm-systems} and
the formative works of \Cref{subsec:formative} were read in full, and the
distinction is drawn here rather than left for a reader to guess.

\section{Existing Solutions and Technologies}
\label{sec:related}

\Cref{sec:literature} surveyed the research. This section surveys what is
actually deployable today, including what the industry partner already runs.

\begingroup
\small
\setlength{\tabcolsep}{5pt}
\begin{xltabular}{\textwidth}{@{}>{\hsize=0.58\hsize}Y>{\hsize=1.06\hsize}Y>{\hsize=1.36\hsize}Y@{}}
  \caption{Deployable technologies in the automated GUI testing space}
  \label{tab:existing}\\
  \toprule
  Category & Representative tools & Strengths and limitations \\
  \midrule
  \endfirsthead
  \caption[]{Deployable technologies in the automated GUI testing space \emph{(continued)}}\\
  \toprule
  Category & Representative tools & Strengths and limitations \\
  \midrule
  \endhead
  \midrule
  \multicolumn{3}{r@{}}{\footnotesize\itshape continued on the next page}\\
  \endfoot
  \bottomrule
  \endlastfoot

  Scripted UI automation frameworks & Appium, Espresso, UI Automator, XCUITest, Selenium, Playwright &
    The industry default. Deterministic, fast, and integrated with continuous integration. They are authoring frameworks, not test designers: every case is written and maintained by a person, and a user-interface redesign invalidates the locators wholesale. They confirm known behaviour and discover nothing. \\
  \addlinespace
  Device farms and execution infrastructure & Firebase Test Lab, AWS Device Farm, BrowserStack, Sauce Labs &
    Solve the device-fragmentation half of the problem: many real devices, on demand. They execute whatever tests they are given and contribute nothing to deciding what to test. Complementary to this project rather than competing with it. \\
  \addlinespace
  Commercial ``AI-assisted'' testing & Testim, mabl, testRigor, Functionize, Applitools &
    Two genuine capabilities: self-healing locators, which re-identify a control after a redesign, and visual regression, which diffs rendered screens. Both reduce maintenance cost on an existing suite. Neither authors new tests from a specification, and the visual approach has no semantic oracle: it detects that a screen changed, not that the change was wrong. \\
  \addlinespace
  Open-source automated explorers & Monkey \cite{monkey}, DroidBot \cite{droidbot}, Humanoid \cite{humanoid}, Stoat, Ape, Fastbot &
    Free, unattended, and effective at finding crashes. Optimise structural coverage; no functional oracle; the output is a crash log rather than a test case a person can read or re-run deliberately. \\
  \addlinespace
  Research agents & The seven systems of \Cref{tab:llm-systems} &
    Demonstrate the capabilities discussed above, but are research artefacts: several provide no replication package (DroidAgent had to re-implement GPTDroid for comparison) and none ingests a specification or defect history. \\
  \addlinespace
  The partner's existing system & The static-knowledge generator described in Section~2 of ETA-TR-2026-001 \cite{etatr} &
    Already ingests an SRS and a Figma export into a Neo4j knowledge graph and runs a multi-stage retrieval-planning loop to generate non-duplicate test cases. Its documented limits are the starting point of this project: retrieval is token-overlap scoring only; execution memory is a verdict and a free-text note per run; and there is no path memory, no defect correlation and no cross-session learning. \\
\end{xltabular}
\endgroup

\subsection*{The gap}

Placing the research and the products side by side yields a clean statement of
the gap. Commercial tools have adoption, reliability and integration but no
test design: they automate execution and maintenance of tests a person wrote.
Research agents have test design but no specification, no durable
cross-application memory, no defect history and no fault attribution: they
explore, and they report crashes. The partner's existing system has the
specification and the graph but no learning: it is a static-knowledge
generator, and a second run knows exactly what the first one knew.

This project addresses the intersection: a system that keeps the partner's
specification grounding, adds the research agents' autonomous execution, and
closes the loop between them by persisting every execution outcome as
structured knowledge that measurably changes the next proposal, while adding
the two things neither side has, namely sound fault attribution and an enforced
safety boundary.

\section{Market and User Research}
\label{sec:market}

\subsection*{Who would use it}

The primary users are quality-assurance engineers in organisations that ship
graphical applications on a short release cadence, the population that today
performs manual exploratory testing, maintains scripted suites, or both. The
secondary users are the application developers who consume the output, whose
requirement is not more reports but fewer false ones. A third group, relevant
to a Bangladeshi context, is the smaller application-development firm that
cannot staff a dedicated QA team at all and currently ships with ad-hoc manual
checks; for them the alternative to an automated exploratory agent is not a
cheaper tool but no systematic testing.

\subsection*{Evidence of demand}

Four independent sources of evidence were available to this project.

\textbf{The industry partner's own commissioning document.} Samsung Research and
Development Institute Bangladesh specified an eight-family enhancement
programme with a stated effort estimate of 17--23 weeks and a phased dependency
ordering~\cite{etatr}. An organisation does not scope four to six months of
engineering against a capability it does not need; the document is the
strongest single demand signal available to the project, and it is specific
about \emph{which} capability, learning across executions, is the valuable
part.

\textbf{Field deployments reported in the literature.} Two of the reviewed
systems were deployed in production organisations rather than only evaluated.
AUITestAgent runs in Meituan's weekly regression process for a video business
line, where its authors report a saving of 1.5 person-days per regression round
and four new functional bugs over two months~\cite{auitestagent}. XUAT-Copilot
has ``launched in the formal testing environment of WeChat Pay mobile app'',
where its authors report that it ``saves a considerable amount of manpower in
the daily development work''~\cite{xuatcopilot}. Both are large payment or
commerce platforms committing production processes to language-model agents,
which is a stronger signal than any published forecast.

\textbf{Evidence about what users want the output to look like.} The developer
study cited by DroidAgent is directly informative for the interface design: the
preferred form for an automatically generated test is \emph{natural language},
not a test-API script, and it should carry expected outputs, the steps of the
use case, and the specific application feature under test~\cite{droidagent}.
This project's generated test cases carry exactly those elements, and the PDF
run report (DR-08) exists because a result that can only be read from a log
does not reach the person who must act on it.

\textbf{Direct engagement with a second application.} A production Bangladeshi
agricultural-commerce application, Shobar Khamar, was supplied by the supervisor as a second
target, with live users and a real payment and one-time-password flow. Its
inclusion was what forced the project's generality requirement (DR-03) and its
safety boundary (DR-07) from aspirations into enforced constraints.

\subsection*{What users do instead, and what it costs them}

The present alternatives are manual exploratory testing, scripted regression
suites, and unattended crash hunting with Monkey-class tools. The first is the
most effective and the least scalable; the second scales but only over
behaviour someone already anticipated; the third is nearly free and finds only
crashes. The measured saving of 1.5 person-days per regression round for a
\emph{single} business line~\cite{auitestagent} gives an order of magnitude for
the manual effort in play.

The evidence in this section is secondary throughout. No structured interview
with the partner's quality-assurance engineers and no practitioner survey was
conducted, so the demand argument rests on published studies and on the
partner's own technical requirements document rather than on primary data
gathered by the team. That is a real limit on how strongly the market claim can
be put, and it is stated here rather than implied.

\section{Feasibility Study}
\label{sec:feasibility}

\subsection*{Technical feasibility}

Feasible, and demonstrated rather than argued: the delivered system executes on
a physical Android device and in a browser, and the individual capabilities
have independent precedent in the literature. The two technical risks are real
and were both encountered. The first is applications whose accessibility tree
is thin: DroidAgent reports the same failure mode, naming applications ``whose
widgets do not have any textual properties'' and single views containing
distinct interactive sub-regions as the cases where its coverage
collapses~\cite{droidagent}, which is exactly the behaviour of the
cross-platform-toolkit application in constraint C4. The second is
non-determinism, which is intrinsic and is managed by evaluating over
distributions of runs rather than single runs (constraint C8).

\subsection*{Economic feasibility}

This is the dimension on which the project's position is strongest, and it
needs stating carefully because the units differ. DroidAgent's cost is reported
per application for a fixed two-hour exploration budget: \$13--\$22, mean
\$18.10~\cite{droidagent}. AutoDroid reports per query: \$10.17--\$18.76 per
thousand GPT-4 queries after optimisation~\cite{autodroid}. This project
measures per \emph{executed test case}, and its figure, reported with its
derivation in \Cref{ch:evaluation}, is of the order of one United States cent,
because roughly nine tenths of each execution's work is performed on the device
rather than in the model, and because the executor and investigator roles were
deliberately moved to cheaper model tiers once the quality cost of doing so had
been measured.

The comparison is therefore indicative rather than like-for-like: a two-hour
DroidAgent campaign and one of this project's test executions are different
units of work. What the comparison does establish is that per-test cost is a
tractable engineering variable rather than a fixed property of the approach,
and that the specific failure DroidAgent reports (being unable to repeat runs
because of API cost, and having to record that as a threat to
validity~\cite{droidagent}) is avoidable. For a student project on a fixed
prepaid balance (constraint C2), that difference is the difference between
running a campaign and not running one.

\subsection*{Operational feasibility}

The system is operable by one engineer: a target is configured as a validated
JSON profile through the dashboard rather than as a code change (DR-09), a
campaign is launched and watched from the same surface, and its result is
exported as a PDF a non-participant can read (DR-08). Two operational
prerequisites remain and are not removable: a physical device or emulator with
the companion accessibility service installed, and disposable test accounts
provisioned in advance for any application behind a sign-in. The deployment
model is the genuine operational boundary: the system is scoped to run as a
local, developer-operated set of processes on a trusted host (constraint C7),
so adoption into an organisation's shared infrastructure would require a
hardening pass that this project did not take on.

\subsection*{Legal and regulatory feasibility}

Feasible under the instruments analysed in \Cref{sec:standards}, subject to the
controls stated there: dedicated disposable test accounts rather than
production user data, destructive and communicative actions blocked at the
configuration layer, application content treated as untrusted data in every
prompt, and credentials kept out of the planner path entirely. No regulatory
obstacle to the project as scoped was identified.

\subsection*{Conclusion and the assumptions it rests on}

The project is feasible on all four dimensions. Its feasibility depends on four
assumptions, each of which held during this work and each of which would need
re-examination for a different target: that the application under test exposes
an accessibility tree or a DOM; that disposable test accounts can be
provisioned in advance for authenticated flows; that a hosted model API remains
available at roughly current prices; and that the organisation accepts a
reported defect as a
\emph{candidate} requiring human confirmation rather than as a verified fault.
The last of these is not a technical assumption but an organisational one, and
on the evidence of \Cref{subsec:measurement} it is the assumption the entire
field currently operates under.
